%!TEX TS-program = pdflatex
%!TEX encoding = UTF-8 Unicode
%!TEX spellcheck = English (Aspell)

\documentclass[11pt]{article}
\usepackage{lucy,handout,graphicx}

% =========================================================
\begin{document}

\headers{CSCI 332}{``Quiz 6'' group activity solutions}{Fall 2026}

\noindent 
For each of the problems below, transform the input into a graph and apply a standard graph algorithm that you’ve seen in class.  Whenever you use a standard graph algorithm, you \EMPH{must} provide the following information.  (I recommend actually using a bulleted list.)
\begin{itemize}\itemsep0pt
\item What are the vertices?  What does each vertex represent?
\item What are the edges?  Are they directed or undirected?
\item If the vertices and/or edges have associated values, what are they?
\item \textbf{What is the precise relationship between your graph and the stated problem?}
\item What standard graph algorithm are you using to solve the problem?
\item What is the running time of your entire algorithm, \emph{including} the time to build the graph, \emph{as a function of the original input parameters}?
\end{itemize}
Finally, it is crucial to remember that even when you are explicitly given a graph as part of the input, that may not be the graph you actually want to search!

\bigskip
\hrule
\medskip

\begin{enumerate}
\parindent 1.5em \itemsep 2ex

%----------------------------------------------------------------------
\item
Describe and analyze an efficient algorithm that either returns the minimum number of moves required to solve a given number maze $M[\range{1}{n}, \range{1}{n}]$, or correctly reports that the maze has no solution.

\begin{Solution}
We construct a directed graph $G=(V,E)$ as follows.
\begin{itemize}
\item
The vertices of $G$ correspond to (indices of) grid cells:
\[
	V := \set{i\mid 1\le i\le n} \times \set{j\mid 1\le i\le n}
\]
There are exactly $n^2$ vertices.

\item
The directed edges of $G$ correspond to legal moves; more formally, $E = E_\leftrightarrow \cup E_\updownarrow$, where
\begin{align*}
	E_\leftrightarrow := {}&
		\Setbar{(i,j)\arcto(i,j’)\strut}{\abs{j-j’} = M[i,j]}
				&& \text{[horizontal edges]} \\
	E_\updownarrow := {}&
		 \Setbar{(i,j)\arcto(i’,j)\strut}{\abs{i-i’} = M[i,j]}
				&& \text{[vertical edges]}
\end{align*}
There are at most $4n^2$ edges.

\item
The vertices and edges do not have associated values.  In particular, each integer $M[i,j]$ is already encoded into the edges leaving vertex $(i,j)$.

\item
Any sequence of legal moves in number maze $M$ corresponds to a directed path in graph $G$, and vice versa.  In particular, the shortest sequence of legal moves from the upper left corner to the lower right corner corresponds to the shortest path in $G$ from $(1,1)$ to $(n,n)$.

\item
We can compute this shortest path using breadth-first search.

\item
Our algorithm runs in $O(V+E) =$ \EMPH{$O(n^2)$ time}.

\end{itemize}
\end{Solution}


%----------------------------------------------------------------------
\newpage
\item
Suppose a sequence of moves is considered \emph{valid} if and only if the moves  alternate between horizontal and vertical.  That is, a valid move sequence never has two horizontal moves in a row or two vertical moves in a row.

Describe and analyze an algorithm that either returns the minimum number of moves in a valid solution to a given number maze, or reports correctly that no valid solution exists.

\begin{Solution}
We construct a directed graph $G=(V,E)$ as follows.

\begin{itemize}
\item
There are two vertices of $G$ for each grid cell, each corresponding to the orientation of the next move out of that cell.
\[
	V := \set{i\mid 1\le i\le n}
			\times \set{j\mid 1\le i\le n}
			\times \set{\updownarrow,\leftrightarrow}
\]
There are exactly $2n^2$ vertices.

\item
The directed edges of $G$ correspond to legal moves; more formally, $E = E_\leftrightarrow \cup E_\updownarrow$, where
\begin{align*}
	E_\leftrightarrow := {}&
		\Setbar{(i,j,\leftrightarrow)\arcto(i,j’,\updownarrow)\strut}
				{\abs{j-j’} = M[i,j]}
				&& \text{[horizontal edges]} \\
	E_\updownarrow := {}&
		 \Setbar{(i,j,\updownarrow)\arcto(i’,j,\leftrightarrow)\strut}
		 		{\abs{i-i’} = M[i,j]}
				&& \text{[vertical edges]}
\end{align*}
There are at most $4n^2$ edges.

\item
The vertices and edges do not have associated values.  In particular, each integer $M[i,j]$ is already encoded into the edges leaving vertex $(i,j)$.

\item
Any sequence of legal moves in number maze $M$ corresponds to a directed path in graph $G$, and vice versa.  In particular, the shortest sequence of legal moves from the upper left corner to the lower right corner corresponds to the shortest path in $G$ from either $(1,1,\updownarrow)$ or $(1,1,\leftrightarrow)$, to either $(n,n,\updownarrow)$ or $(n,n,\leftrightarrow)$.

\item
We can compute all four shortest paths from $(1,1,?)$ to $(n,n,?)$ using two calls two breadth-first search,  starting at $(1,1,\updownarrow)$ and $(1,1,\leftrightarrow)$, and then return the shortest of those four shortest paths.

\item
Our algorithm runs in $O(V+E) =$ \EMPH{$O(n^2)$ time}.

\end{itemize}
The graph $G$ is a standard product construction between the basic number-maze graph, described in our solution to problem 1, and a cycle of length $2$.
\end{Solution}

%----------------------------------------------------------------------
\newpage
\item
Suppose a sequence of moves is considered \emph{valid} if and only if its length (the number of moves) is a multiple of $5$.

Describe and analyze an algorithm that either returns the minimum number of moves in a valid solution to a given number maze, or reports correctly that no valid solution exists.

\begin{Solution}
We construct a directed graph $G=(V,E)$ as follows.
\begin{itemize}
\item
There are five vertices of $G$ for each grid cell, each corresponding to a different remainder mod 5:
\[
	V := \set{i\mid 1\le i\le n}
			\times \set{j\mid 1\le i\le n}
			\times \set{0,1,2,3,4}
\]
There are exactly $5n^2$ vertices.

\item
The directed edges of $G$ correspond to legal moves; more formally, $E = E_\leftrightarrow \cup E_\updownarrow$, where
\begin{align*}
	E_\leftrightarrow := {}&
		\Setbar{(i,j,r)\arcto(i,j’,r+1\bmod 5)\strut}{\abs{j-j’} = M[i,j]}
	\\
	E_\updownarrow := {}&
		 \Setbar{(i,j,r)\arcto(i’,j,r+1\bmod 5)\strut}{\abs{i-i’} = M[i,j]}
\end{align*}
There are at most $20n^2$ edges.

\item
The vertices and edges do not have associated values.  In particular, each integer $M[i,j]$ is already encoded into the edges leaving each vertex $(i,j,r)$.

\item
Any valid sequence of legal moves in number maze $M$ corresponds to a directed path in graph $G$, and vice versa.  In particular, the shortest valid sequence of legal moves from the upper left corner to the lower right corner corresponds to the shortest path in $G$ from $(1,1,0)$ to $(n,n,0)$.

\item
We can compute this shortest path using breadth-first search.

\item
Our algorithm runs in $O(V+E) =$ \EMPH{$O(n^2)$ time}.

\end{itemize}
The graph $G$ is a standard product construction between the basic number-maze graph, described in our solution to problem 1, and a cycle of length $5$.
\end{Solution}
\end{enumerate}


%----------------------------------------------------------------------
\newpage

\bigskip
\noindent
\textbf{Harder problems to think about later:}

\begin{enumerate}
\setcounter{enumi}{3}
\parindent 1.5em \itemsep 2ex

\def\arcTo{\mathord{\Rightarrow}}
\item
Now suppose a sequence of moves is considered \emph{valid} if and only if it contains no U-turns.

Describe and analyze an algorithm that either returns the minimum number of moves in a valid solution to a given number maze, or reports correctly that no valid solution exists.

\begin{Solution}
The main idea is to keep track not only of the position of the token, but the direction of its last move.

First let $G=(V,E)$ be the directed graph of grid squares and legal moves from problem 1.  We construct a new directed graph $G’ = (V’, E’)$ as follows:
\begin{itemize}
\item
The vertices $V’$ correspond to \emph{legal moves} in the number maze, except for a special start vertex indicating that no moves have yet been made.  More formally, we have
\[
	V’ = \set{(1,1)} \cup E.
\]
Each vertex $\arc{u}{v}$ corresponds to the last move taken by the token; vertex $(1,1)$ corresponds to the token’s start in the upper left corner.  There are at most $4n^2 + 1$ vertices in $V’$.

\item
The directed edges $E’$ correspond to legal initial moves and valid consecutive pairs of legal moves:
\begin{itemize}
\item $(1,1)\arcTo((1,1)\arcto v)$ for every legal move $(1,1)\arcto v\in E$
\item $(u\arcto v)\arcTo(v\arcto w)$ for each pair of legal moves $u\arcto v\in E$ and $v\arcto w\in E$ that share a grid cell~$v$ and do not point in opposite directions.
\end{itemize}
(I’m using double arrows to denote directed edges in $E’$, to distinguish them from edges in $E$.)  There are at most $2$ legal moves from the starting cell, and for any legal move there are at most three valid next moves.  Thus, there are at most $2 + 12n^2$ edges in $E’$.

\item
The vertices and edges do not have associated values.

\item
Any valid sequence of $\ell$ legal moves in $M$ corresponds to a directed path in $G’$ with length $\ell$.  For example, the 4-move sequence $(1,1) \arcto (3,1) \arcto (3,7) \arcto (4,7) \arcto (4,4)$ corresponds to the directed path
\[
	(1,1)\arcTo
	((1,1)\arcto(3,1))\arcTo
	((3,1)\arcto(3,7))\arcTo
	((3,7)\arcto(4,7))\arcTo
	((4,7)\arcto(4,4)),
\]
which has length $4$. Thus, the shortest valid sequence of legal moves from $(1,1)$ to $(n,n)$ corresponds to the shortest path in $G$ from the start vertex $(1,1)$ to any vertex of the form $u\arcto(n,n)$.

\item
We can compute this shortest path using breadth-first search.  Specifically, we compute shortest path in $G’$ from $(1,1)$ to every other vertex in $V’$ using breadth-first search, and then find the best last move $u\arcto(n,n)$ by comparing distances by brute force.

\item
Our algorithm runs in $O(V+E) =$ \EMPH{$O(n^2)$ time}.

\end{itemize}
\end{Solution}


%----------------------------------------------------------------------
\newpage
\item
Now suppose a sequence of moves is considered \emph{valid} if and only if the sequence of move \emph{lengths} is increasing.

Describe and analyze an algorithm that either returns the minimum number of moves in a valid solution to a given number maze, or reports correctly that no valid solution exists.

\begin{Solution}
We construct a directed graph $G=(V,E)$ as follows.
\begin{itemize}
\item
The vertices of $G$ correspond to (indices of) grid cells and the length of the move that brought the token to that grid cell:
\[
	V := \set{i\mid 1\le i\le n}
		\times \set{j\mid 1\le i\le n}
		\times \set{\ell\mid 0\le i\le n-1}
\]
There are exactly $n^3$ vertices.

\item
The directed edges of $G$ correspond to legal moves; more formally, $E = E_\leftrightarrow \cup E_\updownarrow$, where
\begin{align*}
	E_\leftrightarrow := {}&
		\Setbar{(i,j,\ell)\arcto(i,j’,\abs{j-j’})\strut}
					{\abs{j-j’} = M[i,j] > \ell}
				&& \text{[horizontal edges]} \\
	E_\updownarrow := {}&
		 \Setbar{(i,j,\ell)\arcto(i’,j,\abs{i-i’})\strut}
		 			{\abs{i-i’} = M[i,j] > \ell}
				&& \text{[vertical edges]}
\end{align*}
There are at most $4n^3$ edges.

\item
The vertices and edges do not have associated values.

\item
Any sequence of legal moves in number maze $M$ corresponds to a directed path in graph $G$, and vice versa.  In particular, the shortest sequence of legal moves from the upper left corner to the lower right corner corresponds to the shortest path in $G$ from the start vertex $(1,1,0)$ to any vertex of the form $(n,n,\ell)$.

\item
We can compute shortest paths from $(1,1,0)$ to every vertex in $V$ using breadth-first search, after which we can find the shortest shortest path to $(n,n,\cdot)$ by brute force.

\item
Our algorithm runs in $O(V+E) =$ \EMPH{$O(n^3)$ time}.

\end{itemize}
This graph $G$ is a \emph{subgraph} of the standard product construction between the basic number-maze graph, described in our solution to problem 1, and a complete directed acyclic graph with $n$ vertices.
\end{Solution}


%----------------------------------------------------------------------
\newpage
\item
Finally, suppose a sequence of moves is considered \emph{valid} if and only if the sequence of move lengths is a palindrome.  (A palindrome is any sequence that is equal to its reversal.)

Describe and analyze an algorithm that either returns the minimum number of moves in a valid solution to a given number maze, or reports correctly that no valid solution exists.

\begin{Solution}
The main idea is to keep track of both ends of the move sequence and move inward.  Let $G=(V,E)$ be the directed graph constructed in our solution ot problem $1$.  We construct a new directed graph $G’=(V’,E’)$ as follows.
\begin{itemize}
\item
The vertices of $G’$ correspond to \emph{pairs} of vertices in $G$
\[
	V’ := V\times V = \Setbar{(i,j;k,l)\strut}{1\le i,j,k,l \le n}
\]
There are exactly $n^4$ vertices.

\item
The directed edges of $G$ correspond to pairs of legal moves, forward from the first cell and backward from the second cell, with the same length
\[
	E’ = \Setbar{(i,j;k,l)\arcto (i’,j’;k’,l’)\strut}
				{\begin{gathered}
					(i,j)\arcto(i’,j’)\in E \\
				 	(k’,l’)\arcto(k,l)\in E \\
					M[i,j] = M[k’,l’]
				\end{gathered}}
\]
There are at most $16n^4$ edges.

\item
The vertices and edges do not have associated values/

\item
Each \emph{palindromic} sequence of legal moves in number maze $M$ corresponds to a directed path in graph $G$, and vice versa.  In particular, the shortest palindromic sequence of legal moves from the upper left corner to the lower right corner corresponds to the shortest path in $G$ from $(1,1;n,n)$ to any vertex of the form $(i,j;i,j)$ (for even-length palindromes) or $(i,j;k,l)$ where $(i,j)\arcto(k,l)\in E$ (for odd-length palindromes).

\item
We can compute shortest paths from $(1,1;n,n)$ to every vertex in $V’$ using breadth-first search, after which we can find the shortest relevant shortest path by brute force comparison.

\item
Our algorithm runs in $O(V+E) =$ \EMPH{$O(n^4)$ time}.

\end{itemize}
The graph $G’$ is a standard product construction between the basic number-maze graph $G$ and its reversal.
\end{Solution}
\end{enumerate}




%----------------------------------------------------------------------


\end{document}